\appendixchapter{Categories and Functors}
\label{app:categories}
\chaptermark{Categories and Functors}
\section*{Introduction}
Category theory is used here as a language for organizing constructions
that are already familiar, not as a separate prerequisite subject. This
appendix introduces only the language used elsewhere in these notes:
categories, functors, contravariance, commutative diagrams, and natural
transformations. The forms examples may be read after Chapter~13; they
organize results proved there and do not supply prerequisites for that
chapter.

Yoneda's lemma, representable functors as a theory, adjoint functors,
general limits and colimits, universal arrows, abelian categories,
exact or derived functors, monads, Kan extensions, and enriched categories
are deliberately outside the scope. The goal is narrow categorical scope
with an uncompressed first encounter.

\section{Objects, Morphisms, and Composition}
\begin{definition}[Category]\label{def:category}
A category $\mathcal C$ consists of objects and, for each ordered pair of
objects $A,B$, a set of \textbf{morphisms}, or arrows,
$\operatorname{Hom}_{\mathcal C}(A,B)$. A morphism written $f:A\to B$ has
source $A$ and target $B$. The category also supplies composition
\[
\operatorname{Hom}(B,C)\times\operatorname{Hom}(A,B)
\longrightarrow\operatorname{Hom}(A,C),\qquad(g,f)\longmapsto g\circ f,
\]
and an identity $1_A:A\to A$ for each object. Whenever the composites are
defined,
\[
h\circ(g\circ f)=(h\circ g)\circ f,
\qquad
1_B\circ f=f=f\circ1_A.
\]
\end{definition}
The arrows are the admissible comparisons or structure-preserving maps of
the subject. They need not be set-theoretic functions.
Objects need not form a set; the familiar categories below may have a
proper class of objects. Their morphisms between any two fixed objects
do form a set. No further size theory is needed here.
\[
\begin{array}{c|l|l}
\text{category}&\text{objects}&\text{morphisms}\\\hline
\mathbf{Set}&\text{sets}&\text{functions}\\
\mathbf{Vect}_{\R}&\text{real vector spaces}&\text{linear maps}\\
\mathbf{Top}&\text{topological spaces}&\text{continuous maps}\\
\mathbf{Man}&\text{smooth manifolds without boundary}&\text{smooth maps}
\end{array}
\]
In each case composition is ordinary map composition; the structural
condition is preserved by composition. Our manifolds are Hausdorff and
second countable as in Chapter~13. Here $\mathbf{Man}$ means manifolds
\emph{without boundary}, for simplicity; Chapter~13's broader treatment
also permits boundary, but that larger category is not needed here.

\begin{example}[Divisibility as a category]
Let $\operatorname{Ob}(\mathcal C)=\N$, and declare that there is a morphism
\[
a\longrightarrow b
\quad\Longleftrightarrow\quad
a\mid b.
\]
Reflexivity $a\mid a$ supplies the identity arrow. If $a\mid b$ and
$b\mid c$, then $a\mid c$, which supplies composition. Between two objects
there is at most one morphism, and that morphism is not being treated as an
ordinary function between the integers. It is an abstract arrow recording
the divisibility relation. For example, $2\mid6$ and $6\mid30$ compose
to the arrow $2\to30$. For three consecutive arrows, either parenthesization
has the same endpoints and must be the unique arrow between them.
This proves associativity; uniqueness likewise proves both identity laws.
\end{example}
More generally, every poset $(P,\leq)$ determines a category with one arrow
$x\to y$ exactly when $x\leq y$. Reflexivity and transitivity are precisely
the identity and composition axioms; antisymmetry says that arrows in both
directions force equality of objects. This links the present language with
Appendix~\ref{app:foundations}. For the poset of subsets of a fixed set,
the arrows are inclusions: $A\subseteq B\subseteq C$ gives the inclusion
$A\to C$. Both orders of composing three inclusions send each element to
itself in the final set. In a general poset the same axioms follow from
uniqueness of arrows with prescribed endpoints.

\begin{definition}[Isomorphism]\label{def:categorical-isomorphism}
A morphism $f:A\to B$ is an isomorphism if there is a morphism $g:B\to A$
with $g\circ f=1_A$ and $f\circ g=1_B$.
\end{definition}
The inverse is unique: if $h$ is another, then
$g=g\circ(f\circ h)=(g\circ f)\circ h=h$.
In the four categories these are bijections, linear isomorphisms,
homeomorphisms (\cref{def:homeomorphism}), and diffeomorphisms.
A bijective continuous map need not be an isomorphism in $\mathbf{Top}$:
the identity from $\R$ with the discrete topology to $\R$ with its usual
topology is continuous and bijective, but its inverse is not continuous.
The category determines which structure an inverse must respect.
Thus equality and isomorphism are different notions. Equal objects are
literally the same object; isomorphic objects may be distinct but have the
same structure as seen by the category. In $\mathbf{Vect}_{\R}$, for
example, two distinct two-dimensional vector spaces are isomorphic after a
basis is chosen, but they are not thereby equal as sets.

\section{Functors and the Direction of a Construction}
\begin{definition}[Covariant functor]\label{def:functor}
A functor $F:\mathcal C\to\mathcal D$ assigns an object $F(A)$ to every
object and a morphism $F(f):F(A)\to F(B)$ to every morphism $f:A\to B$,
with $F(1_A)=1_{F(A)}$ and $F(g\circ f)=F(g)\circ F(f)$.
\end{definition}
The definition has two assignments: an object assignment and an arrow
assignment. Both are essential. Forgetting topology gives
$\mathbf{Top}\to\mathbf{Set}$: it sends a space to its underlying set and a
continuous map to the same underlying function. Forgetting vector operations
similarly gives $\mathbf{Vect}_{\R}\to\mathbf{Set}$.

Another example on $\mathbf{Set}$ is
\[
F(X)=X\times X,
\qquad
F(f)(x_1,x_2)=(f(x_1),f(x_2)).
\]
For an identity function, both coordinates remain unchanged, so
$F(1_X)=1_{F(X)}$. For
$X\xrightarrow{f}Y\xrightarrow{g}Z$, applying first
$F(f)$ and then $F(g)$ sends $(x_1,x_2)$ to
$(g(f(x_1)),g(f(x_2)))$, which is exactly $F(g\circ f)(x_1,x_2)$.
Thus the axioms express compatibility, not decoration. Functors send
isomorphisms to isomorphisms because applying $F$ to the two inverse
equations preserves them.

Another concrete functor sends a set to its power set and a function
$f:X\to Y$ to the direct-image map $A\mapsto f(A)$. The identities
\[
\operatorname{id}_X(A)=A,\qquad (g\circ f)(A)=g(f(A))
\]
verify the functor laws elementwise. Notice that this uses direct image;
inverse image will reverse the arrow direction instead.

\subsection{Contravariance and opposite categories}
Some constructions naturally reverse arrows. A linear map $T:V\to W$
induces pullback of covectors
\[
T^*:W^*\longrightarrow V^*,
\qquad T^*(\lambda)=\lambda\circ T.
\]
The output direction reverses because a covector on the target can be
precomposed with a map from the source. If
$V\xrightarrow{T}W\xrightarrow{S}U$,
then
\[
(S\circ T)^*=T^*\circ S^*.
\]
The order of the two induced arrows is reversed as well.

\begin{definition}[Contravariant functor]
A \textbf{contravariant functor} $F$ from $\mathcal C$ to $\mathcal D$
assigns $F(A)$ to each object and
\[
F(f):F(B)\longrightarrow F(A)
\]
to each $f:A\to B$, with
\[
F(1_A)=1_{F(A)},
\qquad
F(g\circ f)=F(f)\circ F(g).
\]
\end{definition}

Only after seeing the arrow reversal do we package it formally.
\begin{definition}[Opposite category]\label{def:opposite-category}
The opposite category $\mathcal C^{\mathrm{op}}$ has the same objects as
$\mathcal C$ and an arrow $A\to B$ in $\mathcal C^{\mathrm{op}}$ exactly
when $\mathcal C$ has the reversed arrow $B\to A$. Composition is inherited
in the reversed order. A contravariant functor on $\mathcal C$ is therefore
an ordinary covariant functor $\mathcal C^{\mathrm{op}}\to\mathcal D$.
\end{definition}
Thus it assigns $F(f):F(B)\to F(A)$ to $f:A\to B$ and satisfies
$F(g\circ f)=F(f)\circ F(g)$. This is bookkeeping for direction, not
a new operation on the underlying spaces.

For example, the chain $A\xrightarrow{f}B\xrightarrow{g}C$ becomes
\[
C\xrightarrow{g^{\mathrm{op}}}B\xrightarrow{f^{\mathrm{op}}}A,
\qquad f^{\mathrm{op}}\circ g^{\mathrm{op}}=(g\circ f)^{\mathrm{op}}.
\]
These opposite arrows are formal reversals; they do not assert that $f$ or
$g$ has an inverse function.

For fixed $k\geq0$, pullback of forms (\cref{sec:manifold-form-operations})
is the principal example:
\[
\Omega^k:\mathbf{Man}^{\mathrm{op}}\longrightarrow\mathbf{Vect}_{\R},
\qquad
M\longmapsto\Omega^k(M),
\]
\[
f:M\to N
\quad\longmapsto\quad
f^*:\Omega^k(N)\to\Omega^k(M).
\]
In degree zero this is the familiar rule $f^*u=u\circ f$.
In degree one a covector measures a pushed-forward tangent vector:
$(f^*\omega)_p(v)=\omega_{f(p)}(df_pv)$.
Directly,
\[
(\operatorname{id}_M)^*=\operatorname{id}_{\Omega^k(M)},
\qquad
(g\circ f)^*=f^*\circ g^*.
\]
These are exactly the identity and reversed-composition axioms.

\section{Commutative Diagrams and Natural Transformations}
A diagram is \textbf{commutative} when composites along any two directed
paths with the same starting and ending objects agree. The practical reading
rule is: choose the two paths, compose along each in travel order, and compare
the resulting morphisms. For a triangle
\[
\begin{tikzcd}[column sep=large]
A \arrow[rr,"f"] \arrow[dr,"h"'] && B \arrow[dl,"g"]\\
& C &
\end{tikzcd}
\]
commutativity means $g\circ f=h$. For the square
\[
\begin{tikzcd}[column sep=large]
A \arrow[r,"f"] \arrow[d,"u"'] & B \arrow[d,"v"]\\
C \arrow[r,"g"'] & D
\end{tikzcd}
\]
the two routes from $A$ to $D$ give $v\circ f$ and $g\circ u$, so
commutativity asserts $v\circ f=g\circ u$. Diagrams summarize equalities;
they do not replace their proof.

As a concrete triangle, let $i:U\to V$ be a subspace inclusion and
$q:V\to V/U$ the quotient map. The triangle with third arrow the zero
map $U\to V/U$ commutes because $q(i(u))=[u]=[0]$.
For a concrete square, use $f:X\to Y$ on top, $f\times f$ on the bottom,
and the diagonal maps on the two sides. Both paths send $x$ to
$(f(x),f(x))$. Such a square compares a construction across different sets,
which motivates the next definition.

\begin{definition}[Natural transformation]\label{def:natural-transformation}
For functors $F,G:\mathcal C\to\mathcal D$, a natural transformation
$\eta:F\Rightarrow G$ assigns a morphism $\eta_A:F(A)\to G(A)$ to
each object, such that for every $f:A\to B$,
\[
G(f)\circ\eta_A=\eta_B\circ F(f).
\]
\end{definition}
The components must form a compatible family; arbitrary object-by-object
choices are not enough. As an elementary example, the diagonals
$\eta_X:X\to X\times X$, $x\mapsto(x,x)$, give a natural transformation
from the identity functor on sets to the product functor just described.
Both routes send $x$ to $(f(x),f(x))$. The naturality square is
\[
\begin{tikzcd}[column sep=large,row sep=large]
F(A) \arrow[r,"\eta_A"] \arrow[d,"F(f)"'] & G(A) \arrow[d,"G(f)"]\\
F(B) \arrow[r,"\eta_B"'] & G(B).
\end{tikzcd}
\]
Reading its two paths gives exactly
$G(f)\circ\eta_A=\eta_B\circ F(f)$.

For contravariant functors, an arrow $f:M\to N$ induces vertical arrows
from values at $N$ down to values at $M$. The central analytic example is
$d:\Omega^k\Rightarrow\Omega^{k+1}$, with component $d_M$ at $M$.
For a smooth map $f:M\to N$, its naturality square is
\[
\begin{tikzcd}[column sep=large,row sep=large]
\Omega^k(N) \arrow[r,"d_N"] \arrow[d,"f^*"'] &
\Omega^{k+1}(N) \arrow[d,"f^*"]\\
\Omega^k(M) \arrow[r,"d_M"'] & \Omega^{k+1}(M).
\end{tikzcd}
\]
Starting with $\omega\in\Omega^k(N)$, the top-then-right route gives
$f^*(d_N\omega)$, while the left-then-bottom route gives
$d_M(f^*\omega)$. Commutativity means
$f^*(d\omega)=d(f^*\omega)$, the theorem proved in
\cref{prop:manifold-exterior-derivative}. It says that differentiating
forms is compatible with every smooth change of space.

Appendix~\ref{app:cohomology} will define cochain complexes and their
cohomology before using them. Its induced-map construction
(\cref{prop:cochain-induced,prop:cohomology-pullback}) gives another functor,
\[
H_{\mathrm{dR}}^k:\mathbf{Man}^{\mathrm{op}}\to\mathbf{Vect}_{\R}.
\]
Pullback preserves closed forms because $d\omega=0$ implies
$d(f^*\omega)=f^*(d\omega)=0$, and it preserves exact forms because
$\omega=d\eta$ implies $f^*\omega=d(f^*\eta)$. It therefore descends to
de Rham cohomology. This is a forward pointer, not an additional assumption
about cohomology.
Category theory becomes more central in algebraic topology, homological
algebra, and algebraic geometry; the language introduced here is enough
for the comparisons in these notes.

\section*{Exercises}
\addcontentsline{toc}{section}{Exercises}
\begin{exercise}
Verify that a partially ordered set defines a category with one arrow
$a\to b$ exactly when $a\leq b$. Which arrows are isomorphisms?
\end{exercise}
\begin{hexercise}{app-b-2}
Prove the two contravariant functor identities for the algebraic dual
$V\mapsto V^*$. Track domains and codomains in each composition.
\end{hexercise}
\begin{hexercise}{app-b-3}
For the functor $X\mapsto X\times X$, prove that each coordinate projection
is a natural transformation to the identity functor on sets.
\end{hexercise}
